
\subsection{Grammar-proper PAC for finite targets}
\label{subsec:finite-target-proper}

The semantic closure need not have a known finite CFG or
MCFG presentation when the target is infinite.
For a finite target, however, semantic properness
automatically upgrades to \emph{grammar-subclass}
properness and admits a terminating saturation
procedure on an explicit finite positive list.

\begin{theorem}[Finite-target proper grammar PAC]
\label{thm:finite-target-proper-pac}
Fix \(f\ge1\) and \(h:\Sigma^*\to M\), and
let \(1\le d\le f\) and \(q=|h(\Sigma^+)|\).
For the class
\[
 \mathcal S^{\mathrm{fin}}_{f,h}
 :=\{L\subseteq\Sigma^*:L\text{ finite and }
             (f,h)\text{-tuple-substitutable}\},
\]
there exists a single one-sided \emph{proper}
support-\(PAC_{d+1}\) learner, valid also for
original full coordinate slices, which outputs a
finite regular word-language
\[
 \widehat L_m\in\mathcal S^{\mathrm{fin}}_{f,h},
 \qquad \widehat L_m\subseteq L,
\]
with
\[
 m\ge
 \left\lceil\frac{4q^d}{e\eps}\right\rceil
 \left\lceil\log_2\frac1\delta\right\rceil
\]
sufficient for the \((\eps,\delta)\) guarantee.
In particular, for finite targets in the fixed-\(h\)
CFG and fixed-\(h\) MCFG classes, the learner is
proper in those respective \emph{grammar} classes:
every output is a finite regular language, hence
has a CFG and a fan-out-one MCFG presentation.

Given an \emph{explicit finite list} of all positive
word labels seen on the observed slices, the output
grammar is constructible by a terminating
finite-rule saturation algorithm, without knowing
the target size or maximum word length.
This computational claim assumes that the finite
positive list is supplied in complete form; it does
not assert that a membership oracle for infinite
labelled slices reveals the end of that list in
finite time.
\end{theorem}

\begin{proof}
Let \(L\in\mathcal S^{\mathrm{fin}}_{f,h}\)
be the unknown target.  The set \(P_m\) of
positive words extracted from any finite number
of observed coordinate slices is finite:
each \(w\in L\) has only finitely many
arity-\(d\) decompositions \((E,\vec x)\)
with \(E[\vec x]=w\), and \(L\) itself is
finite.  Define
\[
 \widehat L_m:=\operatorname{Cl}_{f,h}(P_m).
\]
By Lemma~\ref{lem:semantic-horn},
\[
 P_m\subseteq\widehat L_m\subseteq L,
 \qquad\widehat L_m\in\mathcal S_{f,h}.
\]
Since \(L\) is finite, \(\widehat L_m\) is finite.
Every finite language is regular and therefore
has a finite CFG and a fan-out-one MCFG
presentation.  This proves both versions of
properness claimed in the theorem.  The PAC
risk bound follows unchanged from
Theorem~\ref{thm:direct-proper-semantic-pac}.

For the terminating algorithm, initialize the
finite working set \(T:=P_m\).  Enumerate every
arity \(1\le r\le f\), every two factor tuples
\(\vec x,\vec y\in(\Sigma^+)^r\) with equal
componentwise \(h\)-type, and contexts
\(E,F\in\mathcal E_r\) for which
\[
 E[\vec x],\ E[\vec y],\ F[\vec y]\in T.
\]
Add every missing conclusion \(F[\vec x]\)
to \(T\), and repeat until a complete pass
adds no new word.

Each pass is finite: all participating words
are finite, and each admits only finitely
many factor/context decompositions.
Every generated word belongs to \(L\),
by induction using the Horn rules and
the target's substitutability.
As \(L\) is finite, at most \(|L|-|P_m|\)
genuinely new words can be added.  Hence
saturation terminates, with final set exactly
\(\operatorname{Cl}_{f,h}(P_m)\).
Output the finite grammar
\[
 S\longrightarrow w
 \qquad(w\in T).
\]
This is a proper finite grammar and uses
only \(f,h\), and the supplied positive list.
\end{proof}

\begin{proposition}[The forced-positive version space]
\label{prop:forced-positive-version-space}
Let \(\mathcal O\) be any realizable
collection of labelled full or support-sensitive
slice observations, and let \(P(\mathcal O)\)
be the set of words occurring at positive
observed coordinates.  Consider all semantic
targets in \(\mathcal S_{f,h}\) consistent
with every positive and negative observed label.
Then
\[
 \bigcap_{\substack{K\in\mathcal S_{f,h}\\
                   K\text{ agrees with }\mathcal O}}
  K
 =\operatorname{Cl}_{f,h}(P(\mathcal O)).
\]
Thus the Horn closure describes \emph{exactly}
the word memberships logically forced by the
observed labelled slices; negative labels
do not force additional positive words.
\end{proposition}

\begin{proof}
Every consistent \(K\) contains \(P(\mathcal O)\),
so by minimality it contains
\(H:=\operatorname{Cl}_{f,h}(P(\mathcal O))\).
Conversely, realizability provides a target
\(L\in\mathcal S_{f,h}\) agreeing with the
observations.  By conservativity
\(H\subseteq L\).  Every observed positive
word lies in \(H\) by construction, and no
observed negative word can lie in \(H\),
since it is absent from \(L\).
Thus \(H\) itself agrees with all observed
labels, belongs to \(\mathcal S_{f,h}\),
and occurs among the \(K\) in the
intersection.  The asserted equality follows.
\end{proof}

\begin{remark}[Finite seeds can have infinite closure]
\label{ex:finite-seed-infinite-closure}
The finite-target assumption above cannot
simply be dropped from the termination proof.
Take \(\Sigma=\{a\}\), trivial typing \(h\),
\(f=1\), and
\[
 P=\{a^2,a^3,a^4\}.
\]
With \(E=a\Box\), the words
\(E[a]=a^2\) and \(E[a^2]=a^3\)
are both in \(P\).
The Horn rule with \(F=\Box\) and \(F[a^2]=a^2\)
forces \(a\).  Reversing the same factor
pair, if \(a^n\) is present for \(n\ge4\)
and \(F=a^{n-1}\Box\), then \(F[a]=a^n\)
forces \(F[a^2]=a^{n+1}\).
Induction yields
\[
 \operatorname{Cl}_{1,h}(P)=a^+.
\]
The reverse inclusion holds because \(a^+\)
is trivially \(h\)-substitutable and
contains \(P\).  Hence a finite seed need
not have finite closure, even when its
closure happens to be regular.

Likewise, the naive finite truncation of
\(a^+\) at length two,
\(\{a,a^2\}\), is \emph{not}
trivially \(h\)-substitutable:
the factors \(a\) and \(a^2\)
share the empty context, but their
distributions differ at \(a\Box\).
Thus grammar-properness for infinite
targets cannot be obtained merely by
truncating the semantic hypothesis.
\end{remark}

\begin{remark}[Exact scope]
The finite-target theorem does not yield
an effective proper learner for
arbitrary infinite fixed-\(h\)
CFG/MCFG targets.  Nor does it collapse
the distinction between PAC anchor
complexity and finite-bit access to
an infinite labelled slice.  It does,
however, give a uniform, target-size
independent sample bound and a concrete
finite CFG/MCFG construction in every
finite-target instance.
\end{remark}
